% Part: computability
% Chapter: computability-theory
% Section: fixed-point-thm

\documentclass[../../../include/open-logic-section]{subfiles}

\begin{document}

\olfileid{cmp}{thy}{fix}
\olsection{قضیهٔ نقطهٔ ثابت}

دوباره مسئلهٔ توقف را در نظر بگیریم. به‌عنوان نمادگذاری موقت،
$\gn{\cfind{x}(y)}$ را برای $\tuple{x,y}$ می‌نویسیم؛ آن را نمایندهٔ
«نامی» برای مقدار $\cfind{x}(y)$ بدانید. با این نمادگذاری می‌توانیم
یکی از برهان‌های تصمیم‌ناپذیری مسئلهٔ توقف را بازگو کنیم.

پرسش: آیا تابع محاسبه‌پذیری چون $h$ با ویژگی زیر وجود دارد؟ به ازای
هر $x$ و~$y$،
\[
h(\gn{\cfind{x}(y)}) =
\begin{cases}
1 & \text{اگر $\cfind{x}(y) \fdefined$} \\
0 & \text{در غیر این صورت.}
\end{cases}
\]

پاسخ: نه؛ وگرنه تابع جزئی
\[
g(x) \simeq
\begin{cases}
0 & \text{اگر $h(\gn{\cfind{x}(x)}) = 0$} \\
\text{تعریف‌نشده} & \text{در غیر این صورت}
\end{cases}
\]
محاسبه‌پذیر می‌بود و بنابراین نمایه‌ای چون~$e$ می‌داشت. اما آنگاه
داریم
\[
\cfind{e}(e) \simeq
\begin{cases}
0 & \text{اگر $h(\gn{\cfind{e}(e)}) = 0$} \\
\text{تعریف‌نشده} & \text{در غیر این صورت،}
\end{cases}
\]
که در این حالت $\cfind{e}(e)$ تعریف شده است اگر و تنها اگر تعریف
نشده باشد؛ تناقض.

اکنون به معادلهٔ دارای $\cfind{e}$ نگاه کنید. در معنایی، نمونه‌ای از
خودارجاعی در آن هست: ترتیب داده‌ایم مقدار $\cfind{e}(e)$ به شیوه‌ای
خاص به $\gn{\cfind{e}(e)}$ وابسته باشد. قضیهٔ نقطهٔ ثابت می‌گوید
به‌طور کلی \emph{می‌توانیم} چنین کنیم---نه‌فقط برای اثبات تناقض‌ها.

\olref{lem:fixed-equiv} دو بیان هم‌ارز از قضیهٔ نقطهٔ ثابت به دست
می‌دهد. از نظر منطقی، هم‌ارزی آن‌ها از درست‌بودن هر دو نتیجه می‌شود؛
اما مقصود واقعی این است که هریک به‌سادگی از دیگری نتیجه می‌شود، پس
می‌توان آن‌ها را بیان‌های جایگزین یک قضیه دانست.

\begin{lem}
\ollabel{lem:fixed-equiv}
گزاره‌های زیر هم‌ارزند:
\begin{enumerate}
\item برای هر تابع محاسبه‌پذیر جزئیِ $g(x,y)$، نمایه‌ای چون~$e$ وجود
  دارد که به ازای هر~$y$،
\[
\cfind{e}(y) \simeq g(e,y).
\]
\item برای هر تابع محاسبه‌پذیر~$f(x)$، نمایه‌ای چون~$e$ وجود دارد که
  به ازای هر~$y$،
\[
\cfind{e}(y) \simeq \cfind{f(e)}(y).
\]
\end{enumerate}
\end{lem}

\begin{proof}
$(1) \Rightarrow (2)$: با داشتن $f$، تابع $g$ را با
$g(x,y)\simeq\fn{Un}(f(x),y)$ تعریف کنید. از (1) نمایه‌ای چون~$e$
بگیرید که به ازای هر~$y$،
\begin{align*}
\cfind{e}(y) & \simeq \fn{Un}(f(e),y) \\
& \simeq \cfind{f(e)}(y).
\end{align*}

$(2) \Rightarrow (1)$: با داشتن $g$، از قضیهٔ $s$-$m$-$n$ تابعی
چون $f$ بگیرید که برای هر $x$ و~$y$،
$\cfind{f(x)}(y)\simeq g(x,y)$. از (2) نمایه‌ای چون~$e$ بگیرید که
\begin{align*}
\cfind{e}(y) & \simeq \cfind{f(e)}(y) \\
& \simeq g(e,y).
\end{align*}
این برهان را کامل می‌کند.
\end{proof}

\begin{explain}
پیش از اثبات گزارهٔ (1) (و در نتیجه (2))، توجه کنید چقدر شگفت است.
$e$ را برنامه‌ای رایانه‌ای بدانید؛ گزارهٔ (1) می‌گوید با داشتن هر
تابع محاسبه‌پذیر جزئیِ $g(x,y)$ می‌توان برنامهٔ رایانه‌ای~$e$ای یافت
که $g_e(y)\simeq g(e,y)$ را محاسبه کند. به بیان دیگر، می‌توان برنامهٔ
رایانه‌ای‌ای یافت که تابعی ارجاع‌کننده به خودِ برنامه را محاسبه کند.
\end{explain}

\begin{thm}
دو گزارهٔ \olref{lem:fixed-equiv} درست‌اند. به‌طور مشخص، برای هر تابع
محاسبه‌پذیر جزئیِ $g(x,y)$، نمایه‌ای چون~$e$ وجود دارد که به ازای
هر~$y$،
\[
\cfind{e}(y) \simeq g(e,y).
\]
\end{thm}

\begin{proof}
اجزای لازم از پیش در بحث مسئلهٔ توقف نهفته‌اند. بگذارید
$\fn{diag}(x)$ تابع محاسبه‌پذیری باشد که به ازای هر~$x$ نمایه‌ای برای
تابع $f_x(y)\simeq\cfind{x}[2](x,y)$ بازمی‌گرداند؛ یعنی
\[
\cfind{\fn{diag}(x)}(y) \simeq \cfind{x}[2](x,y).
\]
$\fn{diag}$ را تابعی بدانید که برنامهٔ تابعی دوموضعی را به برنامهٔ
تابعی یک‌موضعی تبدیل می‌کند؛ تابع اخیر با ثابت‌گرفتن خودِ برنامه به‌عنوان
آرگومان نخست به دست می‌آید. می‌توان $\fn{diag}$ را به‌طور رسمی چنین
تعریف کرد: نخست $s$ را تعریف کنید:
\[
s(x,y) \simeq \fn{Un}^2(x,x,y),
\]
که در آن $\fn{Un}^2$ تابعی سه‌موضعی و جهانی برای توابع محاسبه‌پذیر
جزئیِ دوموضعی است. سپس بنا بر قضیهٔ $s$-$m$-$n$ می‌توان تابع بازگشتی
اولیه‌ای چون $\fn{diag}$ یافت که
\[
\cfind{\fn{diag}(x)}(y) \simeq s(x,y)
\]
را برآورده کند.

اکنون تابع $l$ را چنین تعریف کنید:
\[
l(x,y) \simeq g(\fn{diag}(x),y).
\]
و بگذارید $\gn{l}$ نمایه‌ای برای $l$ باشد. سرانجام بگذارید
$e=\fn{diag}(\gn{l})$. آنگاه برای هر $y$ داریم:
\begin{align*}
\cfind{e}(y) & \simeq \cfind{\fn{diag}(\gn{l})}(y) \\
& \simeq \cfind{\gn{l}}[2](\gn{l}, y) \\
& \simeq l(\gn{l}, y) \\
& \simeq g(\fn{diag}(\gn{l}),y) \\
& \simeq g(e, y),
\end{align*}
چنان‌که لازم بود.
\end{proof}

\begin{explain}
چه اتفاقی می‌افتد؟ فرض کنید وظیفه دارید برنامه‌ای رایانه‌ای بنویسید
که خودش را چاپ کند. افزون بر آن فرض کنید با زبان برنامه‌نویسی‌ای کار
می‌کنید که کتابخانه‌ای غنی و عجیب از توابع رشته‌ای دارد. به‌ویژه، فرض
کنید زبان شما تابعی چون $\fn{diag}$ دارد که چنین کار می‌کند: با گرفتن
رشتهٔ ورودی~$s$، $\fn{diag}$ هر رخداد نماد «x» را در~$s$ پیدا می‌کند
و آن را با نسخه‌ای نقل‌قول‌شده از رشتهٔ اصلی جایگزین می‌کند. برای
نمونه، با رشتهٔ
\begin{quote}
\begin{verbatim}
hello x world
\end{verbatim}
\end{quote}
به‌عنوان ورودی، تابع رشتهٔ
\begin{quote}
\begin{verbatim}
hello 'hello x world' world
\end{verbatim}
\end{quote}
را بیرون می‌دهد. در این حالت نوشتن برنامهٔ مطلوب آسان است؛ می‌توانید
بررسی کنید که
\begin{quote}
\begin{verbatim}
print(diag('print(diag(x))'))
\end{verbatim}
\end{quote}
کار را انجام می‌دهد. همین اندیشه برای زبان‌های رایج‌تری مانند ++C و
جاوا نیز، با پیاده‌سازی‌ای مفصل‌تر، کار می‌کند.

تنها چند گام با برهان قضیهٔ نقطهٔ ثابت فاصله داریم. فرض کنید گونه‌ای
از تابع چاپ، $\fn{print}(x,y)$، رشتهٔ $x$ و آرگومان عددی دیگری چون
$y$ را می‌پذیرد و رشتهٔ $x$ را $y$ بار چاپ می‌کند. آنگاه «برنامهٔ»
\begin{quote}
\begin{verbatim}
getinput(y);
print(diag('getinput(y); print(diag(x), y)'), y)
\end{verbatim}
\end{quote}
روی ورودی $y$، خودش را $y$ بار چاپ می‌کند. جایگزین‌کردن اسکلت
$\fn{getinput}$---$\fn{print}$---$\fn{diag}$ با تابع دلخواه
$g(x,y)$، عبارت
\begin{quote}
\begin{verbatim}
g(diag('g(diag(x), y)'), y)
\end{verbatim}
\end{quote}
را به دست می‌دهد؛ برنامه‌ای که روی ورودی $y$، تابع $g$ را بر خودِ
برنامه و~$y$ اجرا می‌کند. اگر «نقل‌قول‌کردن» را «استفاده از یک نمایه
برای» بدانیم، برهان بالا را داریم.

فعلاً اگر می‌خواهید برهان را بازی صوری یا جادوی سیاه بدانید اشکالی
ندارد. اما باید بتوانید جزئیات استدلال بالا را بازسازی کنید. هنگامی
که قضیه‌های ناتمامیت (و «قضیهٔ نقطهٔ ثابت» مرتبط) را ثابت کنیم،
راه‌های دیگری برای فهم چراییِ کارکرد آن بررسی خواهیم کرد.
\end{explain}

\begin{tagblock}{lambda}
\begin{digress}
از همین اندیشه می‌توان برای به‌دست‌آوردن یک ترکیب‌گر «نقطهٔ ثابت»
استفاده کرد. فرض کنید ترم لامبدای $g$ را دارید و ترم دیگری چون $k$
می‌خواهید که $k$ با $gk$ هم‌ارزِ $\beta$ باشد. با قراردادهای
نمادگذاری ما ترم‌های
\[
\fn{diag}(x) = xx
\]
و
\[
l(x) = g(\fn{diag}(x))
\]
را تعریف کنید؛ به بیان دیگر، $l$ ترم $\lambd[x][g(xx)]$ است. بگذارید
$k$ ترم $ll$ باشد. آنگاه داریم:
\begin{align*}
k & = (\lambd[x][g(xx)])(\lambd[x][g(xx)]) \\
& \red  g((\lambd[x][g(xx)])(\lambd[x][g(xx)])) \\
& = gk.
\end{align*}
اگر بگیریم
\[
Y = \lambd[g][((\lambd[x][g(xx)])(\lambd[x][g(xx)]))]
\]
آنگاه $Yg$ و $g(Yg)$ به ترمی مشترک کاهش می‌یابند؛ پس
$Yg\equiv_\beta g(Yg)$. این «ترکیب‌گر کِری» نام دارد. اگر در عوض
بگیریم
\[
Y = (\lambd[xg][g(xxg)])(\lambd[xg][g(xxg)])
\]
آنگاه در واقع $Yg$ به $g(Yg)$ کاهش می‌یابد، که گزاره‌ای قوی‌تر است.
این نسخهٔ دومِ $Y$ «ترکیب‌گر تورینگ» نام دارد.
\end{digress}
\end{tagblock}

\end{document}
